\documentclass{article}
\input{C:/Users/aryay/OneDrive/Shared Files/Programs/Preamble.tex}


\title{Geodesics on a Cone}
\author{Arya Yae}
\date{}

\begin{document}

\maketitle

\section*{The Geodesic Equation}
%Let $B$ be the unit disk (a 2D ball) centered at the origin in $\C\cong\R^2$.
First, let $g=|\de z|^2=\de x^2+\de y^2$ be the standard metric on $\C\cong\R^2$.  We wish to build a path $\gamma\from[0,\infty)\to\C$ with zero second derivative from the initial data $\gamma(0)=p$ and $\dot\gamma(0)=\vec v$ (the former is the \emph{initial position} and the latter is the \emph{initial velocity} of $\gamma$).  This means we need to solve the differential equation
\begin{equation}
    \ddot\gamma(t)=0
\end{equation}
for all $t$.  It only takes a little work to show the unique solution is $\gamma(t)=p+t\vec v$, which is just the path which goes in a straight line at a constant rate.  This is the prototypical example of a \emph{geodesic}.

We wish to generalize the notion of a geodesic in more interesting geometries.  The trouble is, it's difficult to generalize the meaning of the second derivative of $\gamma$---it fundamentally depends on what coordinate system you're using, and in general there is no one coordinate system that is more natural than any others (think about a sphere or a torus).  We therefore need a new way to track how the velocity vector $\dot\gamma(t)$ changes over time, and what it means for that velocity to be `constant' (again, picture a path going around a sphere; the velocity of $\gamma$ has to change to stay tangent to the sphere, but we can still require that the path does not turn left or right).  The start of Chapter 4 of Lee's \emph{Riemanian Manifolds} describes this much better than I just have.

The correct generalization turns out to be a \emph{connection}, which is closely tied to the \emph{covariant derivative}.  It's a very confusing definition, but what you need to know is that a choice of metric naturally defines a connection (the Levi-Civita connection) which we denote $\nabla$.  Thankfully, in the 2-dimensional case, especially with a conformal metric, the formula is pretty simple: $\nabla=\de+\del\log u^2$ wehre $\del\log u^2=\del_z\log u^2\,\de z$ is a comlpex $1\times 1$ matrix-valued 1-form.  The the geodesic equation is $\nabla_{\dot\gamma}\dot\gamma=0$ (in general, $\nabla_uv$ is obtained by feeding $u$ into the form part and multiplying $v$ by the resulting $1\times1$ matrix).

Say the metric is $g=u^2|\de z|^2$.  The $u=1$ case is the standard metric, and if $u=|z|^{\beta-1}$ we get the conical metric with cone angle $\beta\in(0,1]$ discussed in the M.R.S.\ paper.  It'll be fun to play around with other choices of $u$, too.  In any case, we calculate $\nabla{\dot\gamma}$ first, which is a 1-form (technically, we need to extend $\dot\gamma$ to a vector field on a tubular neighborhood of $\im\gamma$):
\begin{equation}
    \nabla_{\dot\gamma}=\de\dot\gamma+(\del\log u^2)\dot\gamma=\de\dot\gamma+(\del_z\log u^2)\dot\gamma\,\de z.
\end{equation}
Here, we consider $\dot\gamma$ to be a complex number.  Now we evaluate the 1-form on the tangent vector $\dot\gamma$:
\begin{equation}\label{eqn: geodesic equation}
    0=\nabla_{\dot\gamma}\dot\gamma=\ddot\gamma+(\partial_z\log u^2)\dot\gamma\cdot\langle\de z,\dot\gamma\rangle=\ddot\gamma+(\del_z\log u^2)\dot\gamma^2
\end{equation}
or, in $(x,y)$ coordinates with $\dot\gamma=\left(\begin{smallmatrix}\dot\gamma_x\\\dot\gamma_y\end{smallmatrix}\right)$ and recalling $\de z=\de x+i\de y$,
\begin{equation}\label{eqn: geodesic equation in x y}
    0=\ddot\gamma+\underbrace{(\partial_x\log u-i\partial_y\log u)}_{\partial_z\log u^2}(\dot\gamma_x+i\dot\gamma_y)(\dot\gamma_x+i\dot\gamma_x).
\end{equation}
Notice that if $u=1$ then $\log u=0$, so the equation \eqref{eqn: geodesic equation} reduces to $\ddot\gamma=0$.  If $u=|z|^{\beta-1}$ then
    $$\partial_z\log u^2
    %=\frac{\beta-1}2\partial_z\log|z|^2
    =\frac{\beta-1}z$$
(optional: verify this).  We are evaluating this at $z=\gamma(t)$, so the geodesic equation is
\begin{equation}\label{eqn: conical geodesic equation}
    0=\ddot\gamma+\frac{\beta-1}\gamma\dot\gamma^2.
\end{equation}

\section*{Approximating Geodesics on a Cone}
So how do we simulate a geodesic path?  We use \href{https://en.wikipedia.org/wiki/Euler_method}{Euler's method}.  We begin with the initial data $\gamma(0)=p$ and $\dot\gamma(0)=\vec v=\left(\begin{smallmatrix}v_1\\v_2\end{smallmatrix}\right)$.  As a complex number, $\dot\gamma(0)=v_1+iv_2$.  We can rearrange \eqref{eqn: conical geodesic equation} to
\begin{equation}
    \ddot\gamma=\frac{1-\beta}\gamma\dot\gamma^2=\frac{1-\beta}\gamma\left(v_1^2-v_2^2+2iv_1v_2\right).
\end{equation}
Fixing some time increment $\Delta t$, we use the current value of $\dot\gamma$ to update $\gamma$ (by adding $\Delta t\,\dot\gamma$), and then we update $\dot\gamma$ by adding $\Delta t\,\ddot\gamma$ to its current value.  Now we have approximations for $\gamma(\Delta t)$ and $\dot\gamma(\Delta t)$.  We repeat to get $\gamma(2\Delta t)$ and $\dot\gamma(2\Delta t)$, and so on.  Throughout this process, we have to translate between complex numbers and 2D coordinates.  It's nothing too complicated, but it can get confusing.

For other (conformal) metrics, it's the exact same process but we replace $\frac{1-\beta}z$ with whatever $-\partial_z\log u^2$ is, and then evaluate at $z=\gamma$.  I've never attempted this kind of simulation before, but it sounds really fun.  I'd love it if we can create little animations of a beam tracing through a geodesic path.  If you do visualize this on a 2D graph, you need to think of it as looking straight down at a cone, so that it's tip corresponds to the origin $(x,y)=(0,0)$.  If we manage this, we can talk about how to display it as an actual cone in 3D space.

\end{document}